
\documentclass{article}
\usepackage{colortbl}
\usepackage{makecell}
\usepackage{multirow}
\usepackage{supertabular}

\begin{document}

\newcounter{utterance}

\centering \large Interaction Transcript for game `cladder', experiment `full\_v1.5\_default', episode 6216 with Qwen3.5{-}4B.
\vspace{24pt}

{ \footnotesize  \setcounter{utterance}{1}
\setlength{\tabcolsep}{0pt}
\begin{supertabular}{c@{$\;$}|p{.15\linewidth}@{}p{.15\linewidth}p{.15\linewidth}p{.15\linewidth}p{.15\linewidth}p{.15\linewidth}}
    \# & \multicolumn{2}{c}{Player} && \multicolumn{2}{c}{Game Master} \\
    \hline

    \theutterance \stepcounter{utterance}  
    & & & \multicolumn{4}{p{0.6\linewidth}}{
        \cellcolor[rgb]{0.9,0.9,0.9}{
            \makecell[{{p{\linewidth}}}]{
                \texttt{\tiny{[P1$\langle$GM]}}
                \texttt{You are an expert in causal inference. The following question is not a typical commonsense query, but rather a meticulously designed question created by a professor specializing in causal inference, intended to assess the students' mastery of the course content.} \\
\texttt{Imagine a self{-}contained, hypothetical world with only the following conditions, and without any unmentioned factors or causal relationships: Gender has a direct effect on smoking and peanut allergy. Smoking has a direct effect on tar deposit. Tar deposit has a direct effect on peanut allergy. Gender is unobserved. We know that female causes smoking. smoking causes absence of tar deposit. female or high tar deposit causes being allergic to peanuts. We observed the person is non{-}female. Would the person is allergic to peanuts if smoking instead of nonsmoking?} \\
\texttt{Start your answer with 'Yes' or 'No', followed by additional reasoning or evidence to support your explanation.} \\
            }
        }
    }
     \\ \\

    \theutterance \stepcounter{utterance}  
    & \multicolumn{4}{p{0.6\linewidth}}{
        \cellcolor[rgb]{0.561,0.737,0.561}{
            \makecell[{{p{\linewidth}}}]{
                \texttt{\tiny{[P1$\rangle$GM]}}
                \texttt{No} \\
\\ 
\texttt{**Reasoning:**} \\
\\ 
\texttt{To determine the answer, we must construct the causal structure based strictly on the provided conditions and then evaluate the counterfactual query.} \\
\\ 
\texttt{1.  **Construct the Causal Graph:**} \\
\texttt{    *   **Gender (\$G\$):** The prompt states "female causes smoking" and "smoking causes absence of tar deposit." It also states "female or high tar deposit causes being allergic."} \\
\texttt{        *   \$G\_\{female\} \textbackslash{}rightarrow Smoking\$} \\
\texttt{        *   \$Smoking \textbackslash{}rightarrow \textbackslash{}neg Tar\$ (Absence of tar)} \\
\texttt{        *   \$G\_\{female\} \textbackslash{}lor HighTar \textbackslash{}rightarrow Allergy\$} \\
\texttt{    *   **Observed State:** The person is "non{-}female" (\$G = Male\$).} \\
\texttt{    *   **Implications of Non{-}Female:**} \\
\texttt{        *   Since \$G\_\{female\} \textbackslash{}rightarrow Smoking\$, and the person is non{-}female, the direct causal trigger for smoking is absent. In this hypothetical world with "no unmentioned factors," a non{-}female individual does not smoke.} \\
\texttt{        *   Therefore, \$Smoking = False\$.} \\
\texttt{        *   Since \$Smoking \textbackslash{}rightarrow \textbackslash{}neg Tar\$, the absence of smoking means the cause for "absence of tar" is not active. Consequently, the person has **High Tar** (or at least, they do not have the "absence of tar" state).} \\
\texttt{        *   Now check the allergy condition: "female or high tar deposit causes being allergic."} \\
\texttt{        *   The person is non{-}female (False).} \\
\texttt{        *   The person has high tar (True).} \\
\texttt{        *   Therefore, the condition (\$False \textbackslash{}lor True\$) is met. The person **is allergic**.} \\
\\ 
\texttt{2.  **Analyze the Counterfactual Query:**} \\
\texttt{    *   The question asks: "Would the person be allergic to peanuts if **smoking** instead of **nonsmoking**?"} \\
\texttt{    *   This is a counterfactual intervention: \$P(Allergy | do(Smoking))\$.} \\
\texttt{    *   We intervene to set \$Smoking = True\$.} \\
\texttt{    *   **Effect on Tar:** If \$Smoking = True\$, then \$\textbackslash{}neg Tar\$ is True (High Tar becomes False/Absence of tar).} \\
\texttt{    *   **Effect on Allergy:** The condition for allergy is "female OR high tar".} \\
\texttt{        *   The person is still non{-}female (\$G = False\$).} \\
\texttt{        *   With the intervention, the person now has **Absence of Tar** (Low Tar).} \\
\texttt{        *   The condition becomes: (\$False \textbackslash{}lor False\$).} \\
\texttt{        *   The result is **False**.} \\
\\ 
\texttt{3.  **Conclusion:**} \\
\texttt{    *   Under the observed state (nonsmoking), the person is allergic due to high tar deposits.} \\
\texttt{    *   Under the counterfactual state} \\
            }
        }
    }
    & & \\ \\

    \theutterance \stepcounter{utterance}  
    & & & \multicolumn{2}{p{0.3\linewidth}}{
        \cellcolor[rgb]{0.9,0.9,0.9}{
            \makecell[{{p{\linewidth}}}]{
                \texttt{\tiny{[GM$|$GM]}}
                \texttt{No} \\
            }
        }
    }
    & & \\ \\

    \theutterance \stepcounter{utterance}  
    & & & \multicolumn{2}{p{0.3\linewidth}}{
        \cellcolor[rgb]{0.9,0.9,0.9}{
            \makecell[{{p{\linewidth}}}]{
                \texttt{\tiny{[GM$|$GM]}}
                \texttt{no} \\
            }
        }
    }
    & & \\ \\

    \theutterance \stepcounter{utterance}  
    & & & \multicolumn{2}{p{0.3\linewidth}}{
        \cellcolor[rgb]{0.9,0.9,0.9}{
            \makecell[{{p{\linewidth}}}]{
                \texttt{\tiny{[GM$|$GM]}}
                \texttt{game\_result = WIN} \\
            }
        }
    }
    & & \\ \\

\end{supertabular}
}

\end{document}
